% =============================================================================
%  PAES Studio: presentación de repaso (plantilla Beamer 16:9)
%  Lista para Overleaf: Nuevo proyecto → Subir proyecto → compilar con pdfLaTeX.
%
%  Estética minimalista con la paleta Zinc de la app: sin degradados, sin
%  símbolos de navegación, líneas finas y alternativas con su tecla A–E.
%  Comandos:
%    \alternativa{A}{texto}       alternativa
%    \alternativaok{C}{texto}     alternativa correcta (para la diapositiva de solución)
%    \paso{n}{texto}      paso numerado del desarrollo
%    \ejebarra{Eje}{0.42}{42 %}   barra proporcional de frecuencia
%  Entornos: cajaclave, cajatrampa, lineatiempo (filas «fecha & hecho \\»).
%
%  El generador scripts/pipeline/generate_latex_summary.py reemplaza los bloques
%  delimitados por %%<NOMBRE>%% … %%</NOMBRE>%%.
% =============================================================================
\documentclass[aspectratio=169,11pt]{beamer}

\usepackage{lmodern}        % antes de fontenc/babel: no depende de las fuentes EC
\usepackage[T1]{fontenc}
\usepackage[utf8]{inputenc}
\usepackage[spanish,es-noshorthands,es-tabla]{babel}
\spanishplainpercent        % \% sin espaciado automático (evita choques en modo matemático)
\usepackage{textcomp}
\usepackage{amsmath,amssymb}
\usepackage{booktabs}
\usepackage{array}
\usepackage[most]{tcolorbox}

% ── Paleta Zinc ──────────────────────────────────────────────────────────────
\definecolor{zinc950}{HTML}{09090B}
\definecolor{zinc900}{HTML}{18181B}
\definecolor{zinc700}{HTML}{3F3F46}
\definecolor{zinc500}{HTML}{71717A}
\definecolor{zinc300}{HTML}{D4D4D8}
\definecolor{zinc200}{HTML}{E4E4E7}
\definecolor{zinc50}{HTML}{FAFAFA}
\definecolor{esmeralda}{HTML}{047857}
\definecolor{rosa}{HTML}{BE123C}

\newcommand{\sen}{\operatorname{sen}}

% ── Tema mínimo ──────────────────────────────────────────────────────────────
\usefonttheme{professionalfonts}
\setbeamertemplate{navigation symbols}{}
\setbeamercolor{background canvas}{bg=white}
\setbeamercolor{normal text}{fg=zinc900}
\setbeamercolor{structure}{fg=zinc900}
\setbeamercolor{alerted text}{fg=zinc950}
\setbeamercolor{item}{fg=zinc500}
\setbeamerfont{frametitle}{series=\bfseries,size=\large}
\setbeamertemplate{itemize items}{\raisebox{0.15ex}{\color{zinc500}\scriptsize$\bullet$}}
\setbeamertemplate{enumerate items}{\color{zinc500}\insertenumlabel.}
\setbeamersize{text margin left=1.6cm,text margin right=1.6cm}

\makeatletter
\setbeamertemplate{frametitle}{%
  \vspace*{0.9em}%
  {\usebeamerfont{frametitle}\color{zinc950}\insertframetitle}\par
  \ifx\insertframesubtitle\@empty\else{\small\color{zinc500}\insertframesubtitle}\par\fi
  \vspace{0.25em}{\color{zinc200}\rule{\textwidth}{0.4pt}}%
}
\makeatother
\setbeamertemplate{footline}{%
  \hspace*{1.6cm}{\color{zinc500}\tiny\insertshorttitle}\hfill
  {\color{zinc500}\tiny\ttfamily\insertframenumber/\inserttotalframenumber}\hspace*{1.6cm}\par\vspace*{0.7em}%
}
\setbeamertemplate{title page}{%
  \vfill
  {\color{zinc500}\small PAES Studio}\par\vspace{0.6em}
  {\usebeamerfont{title}\bfseries\Huge\color{zinc950}\inserttitle}\par\vspace{0.5em}
  {\large\color{zinc700}\insertsubtitle}\par\vspace{1.2em}
  {\color{zinc300}\rule{0.35\textwidth}{0.4pt}}\par\vspace{0.6em}
  {\small\color{zinc500}\insertdate}\par
  \vfill
}

% ── Componentes ──────────────────────────────────────────────────────────────
\newtcbox{\tecla}{on line, boxrule=0.4pt, arc=1.5pt, colback=zinc50, colframe=zinc300,
  left=3pt, right=3pt, top=1pt, bottom=1pt, fontupper=\ttfamily\footnotesize\color{zinc700}}
\newtcbox{\teclaok}{on line, boxrule=0.4pt, arc=1.5pt, colback=esmeralda, colframe=esmeralda,
  left=3pt, right=3pt, top=1pt, bottom=1pt, fontupper=\ttfamily\footnotesize\bfseries\color{white}}
\newtcolorbox{opcion}[1][]{enhanced, boxrule=0.4pt, arc=2pt, colback=white, colframe=zinc300,
  left=6pt, right=6pt, top=3pt, bottom=3pt, before skip=3pt, after skip=3pt, #1}
\newcommand{\alternativa}[2]{\begin{opcion}\tecla{#1}\hspace{0.6em}#2\end{opcion}}
\newcommand{\alternativaok}[2]{\begin{opcion}[colframe=esmeralda!60, colback=esmeralda!5]\teclaok{#1}\hspace{0.6em}#2\end{opcion}}
\newcommand{\paso}[2]{\par\noindent\makebox[1.4em][l]{\color{zinc500}\ttfamily\small #1}\parbox[t]{\dimexpr\linewidth-1.4em}{#2}\par\vspace{0.45em}}
\newtcolorbox{cajaclave}[1][]{enhanced, boxrule=0.4pt, arc=2pt, colback=zinc50, colframe=zinc300,
  borderline west={2pt}{0pt}{zinc900}, left=7pt, right=7pt, top=3pt, bottom=4pt,
  fonttitle=\bfseries\small, coltitle=zinc900, colbacktitle=zinc50, titlerule=0pt, title={#1}}
\newtcolorbox{cajatrampa}[1][Trampas frecuentes]{enhanced, boxrule=0.4pt, arc=2pt, colback=white, colframe=rosa!35,
  borderline west={2pt}{0pt}{rosa}, left=7pt, right=7pt, top=3pt, bottom=4pt,
  fonttitle=\bfseries\small, coltitle=zinc900, colbacktitle=white, titlerule=0pt, title={#1}}
\newenvironment{lineatiempo}
  {\footnotesize\renewcommand{\arraystretch}{1.15}\begin{tabular}{@{}>{\leavevmode\bfseries\color{zinc700}}p{0.14\textwidth}p{0.74\textwidth}@{}}}
  {\end{tabular}}
\newlength{\anchobarra}
\newcommand{\ejebarra}[3]{%
  \par\noindent\makebox[0.34\linewidth][l]{\small #1}%
  \setlength{\anchobarra}{0.5\linewidth}%
  \makebox[\anchobarra][l]{{\color{zinc200}\rule[0.2ex]{\anchobarra}{1.1ex}}\hspace*{-\anchobarra}{\color{zinc900}\rule[0.2ex]{#2\anchobarra}{1.1ex}}}%
  \hspace{0.8em}{\small\ttfamily #3}\par\vspace{0.5em}}

%%<META>%%
\title[Repaso PAES: Competencia Matemática 1]{Repaso PAES}
\subtitle{Competencia Matemática 1 (M1)}
\date{8 de noviembre de 2026}
%%</META>%%

\begin{document}

\begin{frame}[plain]
  \titlepage
\end{frame}

%%<CONTENIDO>%%
\begin{frame}{Contenidos}

  \framesubtitle{4 ejes, en el orden del temario}

  \begin{enumerate}\small

    \item \textbf{Números}: Números enteros y racionales, Porcentaje, Potencias y raíces

    \item \textbf{Álgebra y funciones}: Expresiones algebraicas, Proporcionalidad directa e inversa, Ecuaciones e inecuaciones lineales, Sistemas de ecuaciones lineales, Función lineal y afín, Función cuadrática

    \item \textbf{Geometría}: Figuras geométricas: Pitágoras, perímetros y áreas, Cuerpos geométricos: áreas y volúmenes, Transformaciones isométricas y vectores, Semejanza y proporcionalidad de figuras

    \item \textbf{Probabilidad y estadística}: Representación de datos: tablas, gráficos y promedio, Medidas de posición: cuartiles, percentiles y diagrama de cajón, Reglas de las probabilidades

  \end{enumerate}

\end{frame}

\begin{frame}{Qué evalúa la prueba}

  \framesubtitle{Distribución por eje en 8 ensayos (435 preguntas)}

  \vspace{0.8em}

  \ejebarra{Números}{0.324}{32 \%}

  \ejebarra{Probabilidad y estadística}{0.232}{23 \%}

  \ejebarra{Álgebra y funciones}{0.228}{23 \%}

  \ejebarra{Geometría}{0.216}{22 \%}

\end{frame}

\section{Números}

\begin{frame}[plain]
  \vfill
  {\color{zinc500}\small Eje}\par\vspace{0.3em}{\Large\bfseries Números}\par\vspace{0.6em}
  {\color{zinc700} Enteros y racionales, porcentaje, potencias y raíces.}\par\vspace{1em}
  {\small\color{zinc500} Números enteros y racionales, Porcentaje, Potencias y raíces}
  \vfill
\end{frame}

\begin{frame}[allowframebreaks]{Números enteros y racionales}
  \framesubtitle{Teoría esencial}
  \begin{itemize}\small
    \item Prioridad: paréntesis, potencias, $\times$ y $\div$, $+$ y $-$.
    \item Restar un negativo es sumar: $a-(-b)=a+b$.
    \item Todo decimal finito o periódico es racional.
    \item Para comparar fracciones, iguala denominadores o pásalas a decimal.
  \end{itemize}
  \begin{cajaclave}[Operatoria con fracciones]\footnotesize
\begin{gather*}
\frac{a}{b}+\frac{c}{d}=\frac{ad+bc}{bd} \\
\frac{a}{b}\cdot\frac{c}{d}=\frac{ac}{bd} \\
\frac{a}{b}:\frac{c}{d}=\frac{a}{b}\cdot\frac{d}{c}
\end{gather*}

Decimal periódico: $0{,}\overline{3}=\dfrac{3}{9}=\dfrac{1}{3}$ y $0{,}1\overline{6}=\dfrac{16-1}{90}=\dfrac{1}{6}$.
  \end{cajaclave}
\end{frame}

\begin{frame}{Desafío}
  \framesubtitle{Números enteros y racionales}
  {\small ¿Cuál es el valor de $-2 + 3\cdot(-4) - (-5)$?}
  \vspace{0.4em}
  \alternativa{A}{$-19$}
  \alternativa{B}{$-9$}
  \alternativa{C}{$1$}
  \alternativa{D}{$9$}
\end{frame}

\begin{frame}[allowframebreaks]{Solución}
  \framesubtitle{Primero la multiplicación; después sumas y restas, cuidando el doble signo.}
  {\small
  \paso{1}{$3\cdot(-4)=-12$.}
  \paso{2}{$-(-5)=+5$.}
  \paso{3}{$-2-12+5=-9$.}
  }
  \alternativaok{B}{$-9$}
  \begin{cajatrampa}\footnotesize
    \begin{itemize}
    \item Sumar numeradores y denominadores: $\frac{1}{2}+\frac{1}{3}\neq\frac{2}{5}$.
    \item Resolver de izquierda a derecha ignorando que la multiplicación va antes que la resta.
    \end{itemize}
  \end{cajatrampa}
\end{frame}

\begin{frame}[allowframebreaks]{Porcentaje}
  \framesubtitle{Teoría esencial}
  \begin{itemize}\small
    \item El 25 \% de 80 es $0{,}25\cdot 80=20$.
    \item Subir 10 \% y luego bajar 10 \% deja el valor un 1 \% más bajo: $1{,}1\cdot 0{,}9=0{,}99$.
    \item La variación porcentual siempre se calcula sobre el valor inicial.
    \item Precio con IVA (19 \%) $=1{,}19\cdot$ precio neto.
  \end{itemize}
  \begin{cajaclave}[Factores de variación]\footnotesize
\begin{gather*}
\text{aumento de } p\,\%:\ x\left(1+\tfrac{p}{100}\right) \\
\text{descuento de } p\,\%:\ x\left(1-\tfrac{p}{100}\right)
\end{gather*}

\[\text{variación porcentual}=\frac{\text{valor final}-\text{valor inicial}}{\text{valor inicial}}\cdot 100\]
  \end{cajaclave}
\end{frame}

\begin{frame}{Desafío}
  \framesubtitle{Porcentaje}
  {\small Un producto cuesta \$12 000 sin IVA. Si el IVA es del 19 \%, ¿cuál es su precio final?}
  \vspace{0.4em}
  \alternativa{A}{\$2280}
  \alternativa{B}{\$12 190}
  \alternativa{C}{\$14 280}
  \alternativa{D}{\$15 600}
\end{frame}

\begin{frame}[allowframebreaks]{Solución}
  \framesubtitle{El precio final es el neto más el 19 \% del neto, es decir, el neto multiplicado por 1,19.}
  {\small
  \paso{1}{IVA: $0{,}19\cdot 12\,000=2280$.}
  \paso{2}{Precio final: $12\,000+2280=14\,280$, o directamente $12\,000\cdot 1{,}19$.}
  }
  \alternativaok{C}{\$14 280}
  \begin{cajatrampa}\footnotesize
    \begin{itemize}
    \item Sumar porcentajes sucesivos: $+20\,\%$ y $-20\,\%$ no se anulan.
    \item Calcular la variación porcentual sobre el valor final.
    \end{itemize}
  \end{cajatrampa}
\end{frame}

\begin{frame}[allowframebreaks]{Potencias y raíces}
  \framesubtitle{Teoría esencial}
  \begin{itemize}\small
    \item Lleva todo a la misma base antes de operar: $8=2^3$, $27=3^3$.
    \item El exponente negativo invierte, no cambia el signo: $2^{-3}=\frac{1}{8}$.
    \item $\sqrt{a+b}\neq\sqrt{a}+\sqrt{b}$.
    \item Notación científica: $0{,}00034=3{,}4\cdot 10^{-4}$.
  \end{itemize}
  \begin{cajaclave}[Propiedades]\footnotesize
\begin{gather*}
a^m\cdot a^n=a^{m+n} \\
\frac{a^m}{a^n}=a^{m-n} \\
(a^m)^n=a^{mn} \\
(ab)^n=a^nb^n
\end{gather*}

\begin{gather*}
\sqrt[n]{ab}=\sqrt[n]{a}\,\sqrt[n]{b} \\
\sqrt[n]{a^m}=a^{m/n} \\
\sqrt{72}=\sqrt{36\cdot 2}=6\sqrt{2}
\end{gather*}
  \end{cajaclave}
\end{frame}

\begin{frame}{Desafío}
  \framesubtitle{Potencias y raíces}
  {\small ¿Cuál es el valor de $\dfrac{9^{2}\cdot 3^{-1}}{27}$?}
  \vspace{0.4em}
  \begin{tcbraster}[raster columns=2, raster column skip=0.6em, raster row skip=0.35em]
  \alternativa{A}{$\dfrac{1}{3}$}
  \alternativa{B}{$1$}
  \alternativa{C}{$3$}
  \alternativa{D}{$9$}
  \end{tcbraster}
\end{frame}

\begin{frame}[allowframebreaks]{Solución}
  \framesubtitle{Todo se puede escribir en base 3.}
  {\small
  \paso{1}{$9^2=(3^2)^2=3^4$ y $27=3^3$.}
  \paso{2}{Numerador: $3^4\cdot 3^{-1}=3^{3}$.}
  \paso{3}{$\dfrac{3^3}{3^3}=1$.}
  }
  \alternativaok{B}{$1$}
  \begin{cajatrampa}\footnotesize
    \begin{itemize}
    \item Multiplicar las bases: $2^3\cdot 2^4\neq 4^7$ (es $2^7$).
    \item Sumar exponentes al elevar una potencia: $(2^3)^2=2^6$, no $2^5$.
    \end{itemize}
  \end{cajatrampa}
\end{frame}

\section{Álgebra y funciones}

\begin{frame}[plain]
  \vfill
  {\color{zinc500}\small Eje}\par\vspace{0.3em}{\Large\bfseries Álgebra y funciones}\par\vspace{0.6em}
  {\color{zinc700} Expresiones algebraicas, proporcionalidad, ecuaciones, sistemas y funciones afín y cuadrática.}\par\vspace{1em}
  {\small\color{zinc500} Expresiones algebraicas, Proporcionalidad directa e inversa, Ecuaciones e inecuaciones lineales, Sistemas de ecuaciones lineales, Función lineal y afín, Función cuadrática}
  \vfill
\end{frame}

\begin{frame}[allowframebreaks]{Expresiones algebraicas}
  \framesubtitle{Teoría esencial}
  \begin{itemize}\small
    \item Primero el factor común: $6x^2-9x=3x(2x-3)$.
    \item Diferencia de cuadrados: $x^2-49=(x+7)(x-7)$.
    \item $(a+b)^2\neq a^2+b^2$: falta el término $2ab$.
  \end{itemize}
  \begin{cajaclave}[Productos notables]\footnotesize
\begin{gather*}
(a\pm b)^2=a^2\pm 2ab+b^2 \\
(a+b)(a-b)=a^2-b^2 \\
(x+p)(x+q)=x^2+(p+q)x+pq
\end{gather*}
  \end{cajaclave}
\end{frame}

\begin{frame}{Desafío}
  \framesubtitle{Expresiones algebraicas}
  {\small ¿Cuál de las siguientes expresiones es equivalente a $4x^2-25$?}
  \vspace{0.4em}
  \alternativa{A}{$(2x-5)^2$}
  \alternativa{B}{$(2x+5)(2x-5)$}
  \alternativa{C}{$(4x-5)(x+5)$}
  \alternativa{D}{$(2x-25)(2x+1)$}
\end{frame}

\begin{frame}[allowframebreaks]{Solución}
  \framesubtitle{Son dos cuadrados perfectos restados: $(2x)^2-5^2$.}
  {\small
  \paso{1}{$4x^2=(2x)^2$ y $25=5^2$.}
  \paso{2}{Diferencia de cuadrados: $(2x)^2-5^2=(2x+5)(2x-5)$.}
  }
  \alternativaok{B}{$(2x+5)(2x-5)$}
  \begin{cajatrampa}\footnotesize
    \begin{itemize}
    \item Olvidar el doble producto en $(a+b)^2$.
    \item Equivocar los signos al sacar un factor negativo: $-x+3=-(x-3)$.
    \end{itemize}
  \end{cajatrampa}
\end{frame}

\begin{frame}[allowframebreaks]{Proporcionalidad directa e inversa}
  \framesubtitle{Teoría esencial}
  \begin{itemize}\small
    \item Pregúntate: si una magnitud aumenta, ¿la otra aumenta o disminuye?
    \item Reparto proporcional: total $\div$ suma de las partes.
    \item Escala 1:50 000: 1 cm en el mapa son 50 000 cm = 500 m reales.
  \end{itemize}
  \begin{cajaclave}[Igualdades clave]\footnotesize
\begin{gather*}
\text{directa: } \frac{y_1}{x_1}=\frac{y_2}{x_2} \\
\text{inversa: } x_1\,y_1=x_2\,y_2
\end{gather*}
  \end{cajaclave}
\begin{center}\footnotesize
\begin{tabular}{@{}>{\raggedright\arraybackslash}p{\dimexpr(\linewidth-4\tabcolsep)/3\relax}>{\raggedright\arraybackslash}p{\dimexpr(\linewidth-4\tabcolsep)/3\relax}>{\raggedright\arraybackslash}p{\dimexpr(\linewidth-4\tabcolsep)/3\relax}@{}}
\toprule
 & \textbf{Directa} & \textbf{Inversa} \\
\midrule
Constante & $\frac{y}{x}=k$ & $x\cdot y=k$ \\
Gráfico & recta que pasa por el origen & hipérbola \\
Ejemplo & kilos comprados y precio & trabajadores y días de obra \\
\bottomrule
\end{tabular}
\end{center}
\end{frame}

\begin{frame}{Desafío}
  \framesubtitle{Proporcionalidad directa e inversa}
  {\small Se reparten \$60 000 entre Ana y Beto en la razón $2:3$. ¿Cuánto recibe Beto?}
  \vspace{0.4em}
  \alternativa{A}{\$24 000}
  \alternativa{B}{\$30 000}
  \alternativa{C}{\$36 000}
  \alternativa{D}{\$40 000}
\end{frame}

\begin{frame}[allowframebreaks]{Solución}
  \framesubtitle{La razón $2:3$ divide el total en 5 partes iguales.}
  {\small
  \paso{1}{Cada parte vale $\dfrac{60\,000}{2+3}=12\,000$.}
  \paso{2}{Beto recibe 3 partes: $3\cdot 12\,000=36\,000$.}
  }
  \alternativaok{C}{\$36 000}
  \begin{cajatrampa}\footnotesize
    \begin{itemize}
    \item Tratar como directa una relación inversa (más trabajadores, menos días).
    \item En un reparto, dividir el total por el número de personas en vez de por la suma de las partes.
    \end{itemize}
  \end{cajatrampa}
\end{frame}

\begin{frame}[allowframebreaks]{Ecuaciones e inecuaciones lineales}
  \framesubtitle{Teoría esencial}
  \begin{itemize}\small
    \item Si $-2x\le 6$, entonces $x\ge -3$ (se invierte la desigualdad).
    \item Para eliminar denominadores, multiplica \textbf{todos} los términos por el m.c.m.
    \item Revisa qué pide la pregunta: a veces no es $x$, sino otra cantidad.
  \end{itemize}
  \begin{cajaclave}[Notación de intervalos (DEMRE)]\footnotesize
\begin{gather*}
[a,b]=\{x: a\le x\le b\} \\
]a,b[=\{x: a<x<b\} \\
[a,+\infty[=\{x: x\ge a\}
\end{gather*}
  \end{cajaclave}
\end{frame}

\begin{frame}{Desafío}
  \framesubtitle{Ecuaciones e inecuaciones lineales}
  {\small La suma de tres números enteros consecutivos es 72. ¿Cuál es el mayor de ellos?}
  \vspace{0.4em}
  \alternativa{A}{23}
  \alternativa{B}{24}
  \alternativa{C}{25}
  \alternativa{D}{26}
\end{frame}

\begin{frame}[allowframebreaks]{Solución}
  \framesubtitle{Llama $x$ al menor; los otros son $x+1$ y $x+2$. Ojo: se pregunta por el mayor.}
  {\small
  \paso{1}{$x+(x+1)+(x+2)=72 \Rightarrow 3x+3=72 \Rightarrow x=23$.}
  \paso{2}{Los números son 23, 24 y 25; el mayor es 25.}
  }
  \alternativaok{C}{25}
  \begin{cajatrampa}\footnotesize
    \begin{itemize}
    \item No invertir la desigualdad al dividir por un número negativo.
    \item Multiplicar por el m.c.m. solo algunos términos.
    \end{itemize}
  \end{cajatrampa}
\end{frame}

\begin{frame}[allowframebreaks]{Sistemas de ecuaciones lineales}
  \framesubtitle{Teoría esencial}
  \begin{itemize}\small
    \item Reducción: elige la incógnita con coeficientes fáciles de igualar.
    \item Comprueba la solución en \textbf{ambas} ecuaciones.
    \item En problemas, define claramente qué representa cada variable.
  \end{itemize}
  \begin{cajaclave}[Cantidad de soluciones de $ax+by=c,\textbackslash\{\} dx+ey=f$]\footnotesize
Única si $\dfrac{a}{d}\neq\dfrac{b}{e}$; ninguna si $\dfrac{a}{d}=\dfrac{b}{e}\neq\dfrac{c}{f}$; infinitas si $\dfrac{a}{d}=\dfrac{b}{e}=\dfrac{c}{f}$.
  \end{cajaclave}
\end{frame}

\begin{frame}{Desafío}
  \framesubtitle{Sistemas de ecuaciones lineales}
  {\small En una librería, 2 cuadernos y 3 lápices cuestan \$5400, y 1 cuaderno y 2 lápices cuestan \$3000. ¿Cuánto cuesta un cuaderno?}
  \vspace{0.4em}
  \alternativa{A}{\$600}
  \alternativa{B}{\$1200}
  \alternativa{C}{\$1800}
  \alternativa{D}{\$2400}
\end{frame}

\begin{frame}[allowframebreaks]{Solución}
  \framesubtitle{Sean $c$ y $\ell$ los precios: $2c+3\ell=5400$ y $c+2\ell=3000$.}
  {\small
  \paso{1}{De la segunda: $c=3000-2\ell$.}
  \paso{2}{Sustituye: $2(3000-2\ell)+3\ell=5400 \Rightarrow 6000-\ell=5400 \Rightarrow \ell=600$.}
  \paso{3}{$c=3000-1200=1800$. Comprobación: $2\cdot 1800+3\cdot 600=5400$.}
  }
  \alternativaok{C}{\$1800}
  \begin{cajatrampa}\footnotesize
    \begin{itemize}
    \item Reemplazar el valor en la misma ecuación de la que se despejó (siempre se cumple).
    \item Responder la variable que no se pedía.
    \end{itemize}
  \end{cajatrampa}
\end{frame}

\begin{frame}[allowframebreaks]{Función lineal y afín}
  \framesubtitle{Teoría esencial}
  \begin{itemize}\small
    \item $m>0$: creciente; $m<0$: decreciente; $m=0$: constante.
    \item Evaluar $f(2)$ es reemplazar $x=2$; resolver $f(x)=2$ es buscar $x$.
    \item Cargo fijo + tarifa por unidad $\Rightarrow f(x)=\text{tarifa}\cdot x+\text{fijo}$.
  \end{itemize}
  \begin{cajaclave}[Pendiente y ecuación]\footnotesize
\begin{gather*}
m=\frac{y_2-y_1}{x_2-x_1} \\
f(x)=m(x-x_1)+y_1
\end{gather*}
  \end{cajaclave}
\end{frame}

\begin{frame}{Desafío}
  \framesubtitle{Función lineal y afín}
  {\small Un taxi cobra \$400 de bajada de bandera y \$150 por cada 200 m recorridos. ¿Qué función entrega el cobro $C$, en pesos, de un viaje de $x$ metros?}
  \vspace{0.4em}
  \alternativa{A}{$C(x)=150x+400$}
  \alternativa{B}{$C(x)=0{,}75x+400$}
  \alternativa{C}{$C(x)=400x+150$}
  \alternativa{D}{$C(x)=200x+150$}
\end{frame}

\begin{frame}[allowframebreaks]{Solución}
  \framesubtitle{La tarifa por metro es $\frac{150}{200}=0{,}75$ pesos; el cargo fijo es 400.}
  {\small
  \paso{1}{Tarifa por metro: $150:200=0{,}75$.}
  \paso{2}{$C(x)=0{,}75x+400$. Control: 200 m cuestan $150+400=550$.}
  }
  \alternativaok{B}{$C(x)=0{,}75x+400$}
  \begin{cajatrampa}\footnotesize
    \begin{itemize}
    \item Restar en distinto orden en el numerador y en el denominador de la pendiente.
    \item Confundir el cargo fijo con la tarifa por unidad.
    \end{itemize}
  \end{cajatrampa}
\end{frame}

\begin{frame}[allowframebreaks]{Función cuadrática}
  \framesubtitle{Teoría esencial}
  \begin{itemize}\small
    \item El signo de $a$ decide si hay máximo o mínimo.
    \item El valor máximo o mínimo es la \textbf{ordenada} del vértice.
    \item Dos puntos simétricos tienen la misma altura: el vértice queda al medio.
  \end{itemize}
  \begin{cajaclave}[Vértice y ceros]\footnotesize
\begin{gather*}
x_v=-\frac{b}{2a},\quad y_v=f(x_v) \\
x=\frac{-b\pm\sqrt{b^2-4ac}}{2a}
\end{gather*}
  \end{cajaclave}
\end{frame}

\begin{frame}{Desafío}
  \framesubtitle{Función cuadrática}
  {\small La altura, en metros, de una pelota lanzada hacia arriba es $h(t)=-5t^2+20t$, con $t$ en segundos. ¿Cuál es la altura máxima?}
  \vspace{0.4em}
  \alternativa{A}{2 m}
  \alternativa{B}{15 m}
  \alternativa{C}{20 m}
  \alternativa{D}{40 m}
\end{frame}

\begin{frame}[allowframebreaks]{Solución}
  \framesubtitle{$a=-5<0$: la parábola tiene un máximo en el vértice.}
  {\small
  \paso{1}{$t_v=-\dfrac{20}{2\cdot(-5)}=2$ s.}
  \paso{2}{$h(2)=-5\cdot 4+20\cdot 2=20$ m.}
  }
  \alternativaok{C}{20 m}
  \begin{cajatrampa}\footnotesize
    \begin{itemize}
    \item Usar $\frac{b}{2a}$ sin el signo menos.
    \item Responder $x_v$ cuando se pide el valor máximo de la función.
    \end{itemize}
  \end{cajatrampa}
\end{frame}

\section{Geometría}

\begin{frame}[plain]
  \vfill
  {\color{zinc500}\small Eje}\par\vspace{0.3em}{\Large\bfseries Geometría}\par\vspace{0.6em}
  {\color{zinc700} Figuras y cuerpos geométricos, transformaciones isométricas y semejanza.}\par\vspace{1em}
  {\small\color{zinc500} Figuras geométricas: Pitágoras, perímetros y áreas, Cuerpos geométricos: áreas y volúmenes, Transformaciones isométricas y vectores, Semejanza y proporcionalidad de figuras}
  \vfill
\end{frame}

\begin{frame}[allowframebreaks]{Figuras geométricas: Pitágoras, perímetros y áreas}
  \framesubtitle{Teoría esencial}
  \begin{itemize}\small
    \item La hipotenusa es el lado mayor del triángulo rectángulo.
    \item Los múltiplos de un trío pitagórico también lo son: (6, 8, 10).
    \item Diámetro $=2\cdot$ radio.
  \end{itemize}
  \begin{cajaclave}[Fórmulas]\footnotesize
Triángulo: $A=\dfrac{b\cdot h}{2}$. Paralelogramo: $A=b\cdot h$. Trapecio: $A=\dfrac{(a+c)\,h}{2}$. Círculo: $A=\pi r^2$, $P=2\pi r$.

\begin{gather*}
a^2+b^2=c^2 \\
\text{tríos: } (3,4,5),\ (5,12,13),\ (8,15,17)
\end{gather*}
  \end{cajaclave}
\end{frame}

\begin{frame}{Desafío}
  \framesubtitle{Figuras geométricas: Pitágoras, perímetros y áreas}
  {\small Un rectángulo mide 9 cm de ancho y 12 cm de largo. ¿Cuánto mide su diagonal?}
  \vspace{0.4em}
  \alternativa{A}{13 cm}
  \alternativa{B}{15 cm}
  \alternativa{C}{21 cm}
  \alternativa{D}{$\sqrt{63}$ cm}
\end{frame}

\begin{frame}[allowframebreaks]{Solución}
  \framesubtitle{La diagonal es la hipotenusa del triángulo rectángulo de catetos 9 y 12.}
  {\small
  \paso{1}{$d^2=9^2+12^2=81+144=225$.}
  \paso{2}{$d=15$ cm (es el trío $3,4,5$ multiplicado por 3).}
  }
  \alternativaok{B}{15 cm}
  \begin{cajatrampa}\footnotesize
    \begin{itemize}
    \item Usar el diámetro como radio.
    \item Confundir área con perímetro.
    \end{itemize}
  \end{cajatrampa}
\end{frame}

\begin{frame}[allowframebreaks]{Cuerpos geométricos: áreas y volúmenes}
  \framesubtitle{Teoría esencial}
  \begin{itemize}\small
    \item $1\ \text{m}^3=1000$ litros y $1\ \text{dm}^3=1$ litro.
    \item Si todas las medidas se multiplican por $k$, el volumen se multiplica por $k^3$.
    \item El manto del cilindro desplegado es un rectángulo de $2\pi r$ por $h$.
  \end{itemize}
  \begin{cajaclave}[Fórmulas]\footnotesize
Paralelepípedo: $V=a\cdot b\cdot c$, $\;A=2(ab+bc+ac)$. Cubo: $V=a^3$, $\;A=6a^2$.

Cilindro: $V=\pi r^2 h$, $\;A=2\pi r^2+2\pi r h$.
  \end{cajaclave}
\end{frame}

\begin{frame}{Desafío}
  \framesubtitle{Cuerpos geométricos: áreas y volúmenes}
  {\small Un estanque cilíndrico tiene 2 m de radio y 1,5 m de altura. Considerando $\pi\approx 3$, ¿cuál es su volumen?}
  \vspace{0.4em}
  \alternativa{A}{$9\ \text{m}^3$}
  \alternativa{B}{$12\ \text{m}^3$}
  \alternativa{C}{$18\ \text{m}^3$}
  \alternativa{D}{$36\ \text{m}^3$}
\end{frame}

\begin{frame}[allowframebreaks]{Solución}
  \framesubtitle{Volumen del cilindro: área de la base por la altura.}
  {\small
  \paso{1}{Base: $\pi r^2\approx 3\cdot 4=12\ \text{m}^2$.}
  \paso{2}{$V\approx 12\cdot 1{,}5=18\ \text{m}^3$.}
  }
  \alternativaok{C}{$18\ \text{m}^3$}
  \begin{cajatrampa}\footnotesize
    \begin{itemize}
    \item Usar el diámetro como radio.
    \item Confundir el área de la superficie con el volumen.
    \end{itemize}
  \end{cajatrampa}
\end{frame}

\begin{frame}[allowframebreaks]{Transformaciones isométricas y vectores}
  \framesubtitle{Teoría esencial}
  \begin{itemize}\small
    \item Reflejar respecto del eje $Y$ cambia el signo de la abscisa.
    \item Una rotación de $180^\circ$ equivale a la simetría respecto del origen.
    \item Vector entre dos puntos: $\overrightarrow{AB}=B-A$.
  \end{itemize}
  \begin{cajaclave}[Reglas en coordenadas]\footnotesize
Traslación: $(x,y)\mapsto(x+a,\,y+b)$. Reflexión eje $X$: $(x,-y)$. Reflexión eje $Y$: $(-x,y)$.

Rotación en $90^\circ$ antihoraria: $(-y,x)$. Rotación en $180^\circ$: $(-x,-y)$.
  \end{cajaclave}
\begin{center}\footnotesize
\begin{tabular}{@{}>{\raggedright\arraybackslash}p{\dimexpr(\linewidth-4\tabcolsep)/3\relax}>{\raggedright\arraybackslash}p{\dimexpr(\linewidth-4\tabcolsep)/3\relax}>{\raggedright\arraybackslash}p{\dimexpr(\linewidth-4\tabcolsep)/3\relax}@{}}
\toprule
\textbf{Transformación} & \textbf{Regla} & \textbf{Conserva} \\
\midrule
Traslación por $(a,b)$ & $(x+a,\,y+b)$ & forma, tamaño y orientación \\
Reflexión eje $Y$ & $(-x,\,y)$ & forma y tamaño \\
Rotación $180^\circ$ & $(-x,\,-y)$ & forma y tamaño \\
\bottomrule
\end{tabular}
\end{center}
\end{frame}

\begin{frame}{Desafío}
  \framesubtitle{Transformaciones isométricas y vectores}
  {\small El punto $P(3,1)$ se rota en $90^\circ$ en sentido antihorario con centro en el origen. ¿Cuáles son las coordenadas de su imagen?}
  \vspace{0.4em}
  \alternativa{A}{$(1,-3)$}
  \alternativa{B}{$(-1,3)$}
  \alternativa{C}{$(-3,-1)$}
  \alternativa{D}{$(3,-1)$}
\end{frame}

\begin{frame}[allowframebreaks]{Solución}
  \framesubtitle{Rotación de $90^\circ$ antihoraria: $(x,y)\mapsto(-y,x)$.}
  {\small
  \paso{1}{$(3,1)\mapsto(-1,3)$.}
  \paso{2}{Comprobación: el punto pasa del primer al segundo cuadrante, como corresponde a un giro de un cuarto de vuelta antihorario.}
  }
  \alternativaok{B}{$(-1,3)$}
  \begin{cajatrampa}\footnotesize
    \begin{itemize}
    \item Confundir la reflexión en el eje $X$ con la del eje $Y$.
    \item Restar el vector de traslación en vez de sumarlo.
    \end{itemize}
  \end{cajatrampa}
\end{frame}

\begin{frame}[allowframebreaks]{Semejanza y proporcionalidad de figuras}
  \framesubtitle{Teoría esencial}
  \begin{itemize}\small
    \item Sombras a la misma hora $\Rightarrow$ triángulos semejantes.
    \item Si las longitudes se duplican, el área se cuadruplica.
    \item Empareja los lados opuestos a ángulos iguales.
  \end{itemize}
  \begin{cajaclave}[Relaciones clave]\footnotesize
\begin{gather*}
\frac{a'}{a}=\frac{b'}{b}=k \\
\text{perímetros en razón } k \\
\text{áreas en razón } k^2
\end{gather*}

Escala $1:n$: 1 cm del plano equivale a $n$ cm de la realidad.
  \end{cajaclave}
\end{frame}

\begin{frame}{Desafío}
  \framesubtitle{Semejanza y proporcionalidad de figuras}
  {\small Un poste de 6 m proyecta una sombra de 4 m. A la misma hora, un árbol proyecta una sombra de 10 m. ¿Cuál es la altura del árbol?}
  \vspace{0.4em}
  \alternativa{A}{12 m}
  \alternativa{B}{15 m}
  \alternativa{C}{16 m}
  \alternativa{D}{20 m}
\end{frame}

\begin{frame}[allowframebreaks]{Solución}
  \framesubtitle{A la misma hora, los triángulos que forman objeto y sombra son semejantes.}
  {\small
  \paso{1}{$\dfrac{h}{10}=\dfrac{6}{4}$.}
  \paso{2}{$h=10\cdot 1{,}5=15$ m.}
  }
  \alternativaok{B}{15 m}
  \begin{cajatrampa}\footnotesize
    \begin{itemize}
    \item Emparejar lados que no se corresponden.
    \item Usar la razón $k$ para áreas, que están en razón $k^2$.
    \end{itemize}
  \end{cajatrampa}
\end{frame}

\section{Probabilidad y estadística}

\begin{frame}[plain]
  \vfill
  {\color{zinc500}\small Eje}\par\vspace{0.3em}{\Large\bfseries Probabilidad y estadística}\par\vspace{0.6em}
  {\color{zinc700} Representación de datos, medidas de posición y reglas de las probabilidades.}\par\vspace{1em}
  {\small\color{zinc500} Representación de datos: tablas, gráficos y promedio, Medidas de posición: cuartiles, percentiles y diagrama de cajón, Reglas de las probabilidades}
  \vfill
\end{frame}

\begin{frame}[allowframebreaks]{Representación de datos: tablas, gráficos y promedio}
  \framesubtitle{Teoría esencial}
  \begin{itemize}\small
    \item Para la mediana, \textbf{ordena} primero los datos.
    \item Con $n$ par, la mediana es el promedio de los dos datos centrales.
    \item Un dato extremo mueve la media, pero casi no mueve la mediana.
  \end{itemize}
  \begin{cajaclave}[Fórmulas]\footnotesize
\begin{gather*}
\bar{x}=\frac{\sum x_i\,f_i}{n} \\
\text{posición de la mediana}=\frac{n+1}{2} \\
f_r=\frac{f}{n}
\end{gather*}
  \end{cajaclave}
\end{frame}

\begin{frame}{Desafío}
  \framesubtitle{Representación de datos: tablas, gráficos y promedio}
  {\small Las notas de un curso se resumen en la tabla:

\begin{center}\small
\begin{tabular*}{\linewidth}{@{\extracolsep{\fill}}lllll@{}}
\toprule
\textbf{Nota} & \textbf{4} & \textbf{5} & \textbf{6} & \textbf{7} \\
\midrule
Estudiantes & 3 & 5 & 8 & 4 \\
\bottomrule
\end{tabular*}
\end{center}

¿Cuál es la mediana de las notas?}
  \vspace{0.4em}
  \begin{tcbraster}[raster columns=2, raster column skip=0.6em, raster row skip=0.35em]
  \alternativa{A}{5}
  \alternativa{B}{5,65}
  \alternativa{C}{6}
  \alternativa{D}{6,5}
  \end{tcbraster}
\end{frame}

\begin{frame}[allowframebreaks]{Solución}
  \framesubtitle{Hay 20 datos: la mediana es el promedio de los datos en las posiciones 10 y 11.}
  {\small
  \paso{1}{Frecuencias acumuladas: 3, 8, 16, 20.}
  \paso{2}{Las posiciones 10 y 11 caen en la nota 6 (posiciones 9 a 16).}
  \paso{3}{Mediana $=\dfrac{6+6}{2}=6$. (La media sería $\frac{113}{20}=5{,}65$.)}
  }
  \alternativaok{C}{6}
  \begin{cajatrampa}\footnotesize
    \begin{itemize}
    \item Calcular la mediana sin ordenar los datos.
    \item Promediar los valores de la tabla sin considerar sus frecuencias.
    \end{itemize}
  \end{cajatrampa}
\end{frame}

\begin{frame}[allowframebreaks]{Medidas de posición: cuartiles, percentiles y diagrama de cajón}
  \framesubtitle{Teoría esencial}
  \begin{itemize}\small
    \item Percentil alto no significa nota alta, sino posición relativa alta.
    \item En un diagrama de caja, la caja contiene el 50 \% central de los datos.
    \item Un tramo más largo del cajón no tiene más datos: tiene los datos más separados.
  \end{itemize}
  \begin{cajaclave}[Definiciones]\footnotesize
$Q_1=P_{25}$, $\;Q_2=P_{50}=$ mediana, $\;Q_3=P_{75}$. Percentil 80: el 80 \% de los datos es menor o igual que él.
  \end{cajaclave}
\end{frame}

\begin{frame}{Desafío}
  \framesubtitle{Medidas de posición: cuartiles, percentiles y diagrama de cajón}
  {\small Ana obtuvo un puntaje ubicado en el percentil 80 de su curso. ¿Qué significa este resultado?}
  \vspace{0.4em}
  \alternativa{A}{Ana respondió correctamente el 80 \% de la prueba.}
  \alternativa{B}{Ana obtuvo 80 puntos.}
  \alternativa{C}{Aproximadamente el 80 \% del curso obtuvo un puntaje menor o igual que el de Ana.}
  \alternativa{D}{Solo el 20 \% del curso aprobó la prueba.}
\end{frame}

\begin{frame}[allowframebreaks]{Solución}
  \framesubtitle{El percentil describe la posición de un dato respecto de los demás, no la cantidad de aciertos.}
  {\small
  \paso{1}{Percentil 80: el 80 \% de los datos es menor o igual que ese valor.}
  \paso{2}{Por eso la alternativa C describe el resultado; A y B confunden posición con puntaje.}
  }
  \alternativaok{C}{Aproximadamente el 80 \% del curso obtuvo un puntaje menor o igual que el de Ana.}
  \begin{cajatrampa}\footnotesize
    \begin{itemize}
    \item Interpretar el percentil 80 como «obtuvo el 80 \% del puntaje».
    \item Leer el largo de la caja como el rango total.
    \end{itemize}
  \end{cajatrampa}
\end{frame}

\begin{frame}[allowframebreaks]{Reglas de las probabilidades}
  \framesubtitle{Teoría esencial}
  \begin{itemize}\small
    \item «Al menos uno» $=1-P(\text{ninguno})$.
    \item Con dos dados hay $6\cdot 6=36$ resultados ordenados.
    \item Eventos independientes: se multiplican sus probabilidades.
  \end{itemize}
  \begin{cajaclave}[Reglas]\footnotesize
\begin{gather*}
P(A)=\frac{\text{casos favorables}}{\text{casos posibles}} \\
P(A^c)=1-P(A)
\end{gather*}

\begin{gather*}
P(A\cup B)=P(A)+P(B)-P(A\cap B) \\
\text{independientes: } P(A\cap B)=P(A)\cdot P(B)
\end{gather*}
  \end{cajaclave}
\end{frame}

\begin{frame}{Desafío}
  \framesubtitle{Reglas de las probabilidades}
  {\small Se lanza una moneda dos veces. ¿Cuál es la probabilidad de obtener al menos una cara?}
  \vspace{0.4em}
  \begin{tcbraster}[raster columns=2, raster column skip=0.6em, raster row skip=0.35em]
  \alternativa{A}{$\dfrac{1}{4}$}
  \alternativa{B}{$\dfrac{1}{2}$}
  \alternativa{C}{$\dfrac{2}{3}$}
  \alternativa{D}{$\dfrac{3}{4}$}
  \end{tcbraster}
\end{frame}

\begin{frame}[allowframebreaks]{Solución}
  \framesubtitle{Usa el complemento: «al menos una cara» es lo contrario de «ninguna cara».}
  {\small
  \paso{1}{$P(\text{ninguna cara})=P(\text{SS})=\dfrac{1}{4}$.}
  \paso{2}{$P(\text{al menos una})=1-\dfrac{1}{4}=\dfrac{3}{4}$.}
  }
  \alternativaok{D}{$\dfrac{3}{4}$}
  \begin{cajatrampa}\footnotesize
    \begin{itemize}
    \item Contar pares no ordenados en experimentos con orden, como dos dados.
    \item Suponer equiprobables resultados que no lo son, como las sumas de dos dados.
    \end{itemize}
  \end{cajatrampa}
\end{frame}
%%</CONTENIDO>%%

\end{document}
